\section{A three-variable ANOVA calculation}
\label{sec:book-three-variable-anova-primer}

The relation between a dynamic edge and a higher-order interaction can be calculated with two
binary source coordinates.  Let $x_j,x_k\in\{-1,+1\}$ and let
$A,B,C\in\mathcal H_i^0$ be three target-contrast vectors.  Consider the target-centered score
\begin{equation}
\Phi(x_j,x_k)
=\phi_0+A x_j+B x_k+C x_jx_k,
\qquad \phi_0\in\mathcal H_i^0.
\label{eq:book-three-variable-score}
\end{equation}
Under the uniform product reference, the four scalar functions
$1,x_j,x_k,x_jx_k$ are orthonormal.  They respectively encode a target-only term, two pair terms,
and one target--source--source term.

Audit the response to source $j$ while holding $x_k$ fixed.  Since the centered binary source
space is one-dimensional, identify the edge operator with its coefficient multiplying $x_j$.
Equation~\ref{eq:book-three-variable-score} then gives
\begin{equation}
G_{i\leftarrow j}(x_k)=A+C x_k.
\label{eq:book-three-variable-response}
\end{equation}
The $A$ direction persists across backgrounds.  The $C$ direction changes sign when the second
source changes.  With $X_k$ uniform,
\begin{equation}
\overline G_{i\leftarrow j}
=\mathbb E_{X_k}G_{i\leftarrow j}(X_k)=A,
\qquad
\delta G_{i\leftarrow j}(X_k)=C X_k,
\label{eq:book-three-variable-mean-residual}
\end{equation}
and therefore
\begin{equation}
\mathbb E_{X_k}
\|\delta G_{i\leftarrow j}(X_k)\|_{p_i}^2
=\|C\|_{p_i}^2.
\label{eq:book-three-variable-dynamic-variance}
\end{equation}
The dynamic variance is exactly the squared mass of the component containing the audited source
$j$ and the additional source $k$.  The general interaction-hierarchy theorem replaces these two
binary coordinates by arbitrary categorical source sets and sums the corresponding orthogonal
masses.

The word ``exactly'' uses the product reference.  To see the role of that geometry, keep balanced
binary marginals but let a correlated law $\mu$ satisfy
\begin{equation}
\mathbb E_\mu[X_jX_k]=\rho\ne0.
\end{equation}
Then
\begin{equation}
\langle x_j,x_k\rangle_{L^2(\mu)}=\rho,
\qquad
\langle 1,x_jx_k\rangle_{L^2(\mu)}=\rho,
\label{eq:book-three-variable-correlated-inner-products}
\end{equation}
and, for example,
\begin{equation}
\mathbb E_\mu\|A X_j+B X_k\|_{p_i}^2
=\|A\|_{p_i}^2+\|B\|_{p_i}^2
+2\rho\langle A,B\rangle_{p_i}.
\label{eq:book-three-variable-correlated-cross-term}
\end{equation}
The algebraic section $A+C x_k$ still describes the finite response at a fixed background.
Its attribution to a sum of mutually orthogonal interaction orders now requires a new conditional
geometry or must retain the cross terms induced by $\mu$.
